\section{Coherencia con el proceso de búsqueda}
\label{sec:coherencia_busqueda}

% Verificar que las features evolucionan a lo largo del trial de
% manera consistente con la dinámica esperada del proceso bayesiano.
%
% Ejemplos de patrones esperados:
%   - entropia_posterior debería decrecer monótonamente (en promedio)
%     a medida que el modelo acumula evidencia.
%   - ganancia_informacion debería ser positiva y mayor en las
%     primeras fijaciones.
%   - eig_optimo debería disminuir a medida que la incertidumbre
%     se reduce.
%
% Análisis de fijaciones consecutivas en la misma celda:
%   - Frecuencia con que el sujeto refija la misma celda.
%   - Impacto en las features dinámicas (e.g. evidencia_visual
%     no cambia si la celda es la misma).
%
% Incluir visualizaciones de la evolución media de cada feature
% en función del índice de fijación t.
